% ===================================================================
%  Семинар 5. Содержательная часть — единый источник для обоих файлов.
%
%  Состав и порядок задач — по «Семинар 5.docx»:
%      Задание 1 — смещение и состоятельность оценки;
%      Задание 2 — мультиколлинеарность по вспомогательным R^2 (4 варианта);
%      Задание 3 — способы борьбы с мультиколлинеарностью, R^2_adj, AIC, BIC;
%      Задание 4 — парная регрессия национального дохода на денежную массу
%                  («(См. эксель)», данные — «Семинар 5 (приложение).xlsx»).
%  Формулы из «Семинар 5.docx» (VIF с порогом 5, R^2, R^2_adj, AIC, BIC,
%  ст. ошибка через остаточную сумму вспомогательной регрессии) вынесены
%  в раздел «Основные формулы».
%
%  Источник задания 4: Магнус Я. Р., Катышев П. К., Пересецкий А. А.
%  Эконометрика. Начальный курс, упр. 2.22 (табл. 2.4). Пункты а), б) —
%  как в «Семинар 5.docx»; пункт в) — из исходного условия учебника.
%
%  ОБОЗНАЧЕНИЯ. Как во всех семинарах: RSS — объяснённая сумма квадратов,
%  ESS — остаточная. В «Семинар 5.docx» в формулах AIC/BIC и ст. ошибки
%  остаточная сумма обозначена RSS; здесь она переписана как ESS.
%
%  Раздел «Основные формулы» и всё внутри \begin{solution}...\end{solution}
%  попадает только в seminar_05_solutions.pdf. Блоки \begin{excel}...\end{excel}
%  печатаются в обоих файлах.
%
%  Числа задания 4 пересчитаны независимо и совпадают с выводом Пакета
%  анализа в «Семинар 5 (приложение).xlsx».
% ===================================================================

\seminartitle{Семинар 5}{Свойства оценок, мультиколлинеарность\\ и выбор модели}

% ===================================================================
\ifsolutions
\section{Основные формулы}

% -------------------------------------------------------------------
\subsection{Свойства оценок}\label{sec:props}

Пусть $\widehat{\beta}_n$ — оценка параметра $\beta$ по выборке объёма $n$.
\begin{equation}
  \operatorname{bias}(\widehat{\beta}_n) = \E\widehat{\beta}_n - \beta,
  \qquad
  \operatorname*{plim}_{n\to\infty}\widehat{\beta}_n = \beta
  \iff
  \lim_{n\to\infty}\Pr\bigl\{|\widehat{\beta}_n - \beta| > \varepsilon\bigr\} = 0
  \ \ \forall\varepsilon > 0 .
  \label{eq:bias}
\end{equation}
Оценка \emph{несмещённая}, если смещение равно нулю при любом $n$;
\emph{состоятельная}, если она сходится по вероятности к $\beta$.
Достаточное условие состоятельности: $\E\widehat{\beta}_n \to \beta$ и
$\D(\widehat{\beta}_n) \to 0$ при $n \to \infty$.

% -------------------------------------------------------------------
\subsection{Мультиколлинеарность и VIF}\label{sec:vif}

\emph{Мультиколлинеарность} — тесная линейная связь между регрессорами.
Диагностика — вспомогательная регрессия $j$-го регрессора на все остальные
(с константой), её коэффициент детерминации $R_j^2$:
\begin{equation}
  \mathrm{VIF}_j = \frac{1}{1 - R_j^2}.
  \label{eq:vif}
\end{equation}
Вывод делается и для каждого коэффициента, и для модели в целом: при
$\mathrm{VIF}_j > 5$ мультиколлинеарность заметная, при $\mathrm{VIF}_j > 10$
($R_j^2 > 0{,}9$) — серьёзная.

\textbf{Стандартная ошибка через вспомогательную регрессию.} Пусть
$\ESS_j$ — остаточная сумма квадратов вспомогательной регрессии $x_j$ на
остальные регрессоры (с константой). Тогда
\begin{equation}
  s_{\widehat{\beta}_j} = \sqrt{\widehat{\D}(\widehat{\beta}_j)}
  = \frac{\widehat{\sigma}}{\sqrt{\ESS_j}},
  \qquad
  \ESS_j = (1 - R_j^2)\sum_t(x_{tj} - \overline{x}_j)^2,
  \label{eq:se}
\end{equation}
поэтому
$s^2_{\widehat{\beta}_j} = \mathrm{VIF}_j\cdot\widehat{\sigma}^2/\sum_t(x_{tj} - \overline{x}_j)^2$:
чем сильнее $x_j$ объясняется остальными регрессорами, тем меньше у него
«собственной» вариации $\ESS_j$ и тем больше ошибка коэффициента. В парной
регрессии вспомогательная регрессия — это регрессия $x$ на одну константу, и
$\ESS_x = \sum(x_t - \overline{x})^2$.

% -------------------------------------------------------------------
\subsection{Качество подгонки и выбор модели}\label{sec:fit}

$\TSS = \sum(y_t - \overline{y})^2$, $\RSS = \sum(\widehat{y}_t - \overline{y})^2$
— объяснённая, $\ESS = \sum(y_t - \widehat{y}_t)^2$ — остаточная сумма
квадратов, $n$ — число наблюдений, $k$ — число параметров вместе с
константой.
\begin{gather}
  R^2 = \frac{\RSS}{\TSS} = 1 - \frac{\ESS}{\TSS},
  \qquad
  R^2_{\text{adj}} = 1 - \frac{\ESS/(n-k)}{\TSS/(n-1)}
                   = 1 - \frac{n-1}{n-k}\bigl(1 - R^2\bigr),
  \label{eq:r2} \\[2pt]
  \mathrm{AIC} = \frac{2k}{n} + \ln\frac{\ESS}{n},
  \qquad
  \mathrm{BIC} = \frac{k\ln n}{n} + \ln\frac{\ESS}{n}.
  \label{eq:ic}
\end{gather}

% -------------------------------------------------------------------
\subsection{Парная регрессия}\label{sec:pair}

Модель $Y_t = \alpha + \beta X_t + \varepsilon_t$; $x_t = X_t - \overline{X}$,
$y_t = Y_t - \overline{Y}$.
\begin{gather}
  \widehat{\beta} = \frac{\sum x_ty_t}{\sum x_t^2},
  \qquad
  \widehat{\alpha} = \overline{Y} - \widehat{\beta}\,\overline{X},
  \qquad
  \RSS = \widehat{\beta}^2\sum x_t^2,
  \qquad
  \ESS = \sum y_t^2 - \RSS,
  \label{eq:ols} \\[2pt]
  s^2 = \frac{\ESS}{n-2},
  \qquad
  s^2_{\widehat{\beta}} = \frac{s^2}{\sum x_t^2},
  \qquad
  s^2_{\widehat{\alpha}} = s^2\Bigl(\frac{1}{n} + \frac{\overline{X}^2}{\sum x_t^2}\Bigr),
  \label{eq:s2} \\[2pt]
  t = \frac{\widehat{\beta} - \beta^0}{s_{\widehat{\beta}}} \sim t(n-2),
  \qquad
  \widehat{\beta} \pm t_{0{,}025}(n-2)\,s_{\widehat{\beta}} .
  \label{eq:t}
\end{gather}
\fi

% ===================================================================
\section{Задачи}

% -------------------------------------------------------------------
\problem{1}

Напишите величину смещения смещённой оценки $\widehat{\beta}$ параметра
$\beta$. Какая оценка $\widehat{\beta}$ параметра $\beta$ называется
состоятельной?

\begin{solution}
\textbf{Смещение} — разность между математическим ожиданием оценки и истинным
значением параметра:
\[
  \operatorname{bias}(\widehat{\beta}) = \E\widehat{\beta} - \beta .
\]
Оценка несмещённая, если $\E\widehat{\beta} = \beta$: в среднем по многим
выборкам она попадает в истинное значение.

\textbf{Состоятельность.} Оценка $\widehat{\beta}_n$ состоятельна, если она
сходится по вероятности к истинному значению \eqref{eq:bias}:
\[
  \lim_{n\to\infty}\Pr\bigl\{|\widehat{\beta}_n - \beta| > \varepsilon\bigr\} = 0
  \quad\text{для любого } \varepsilon > 0 .
\]
По мере роста выборки оценка всё ближе к истинному значению: вероятность
ошибиться больше чем на любое заданное $\varepsilon$ стремится к нулю. Обычно
состоятельность требуют в дополнение к несмещённости.

\begin{remark}
Несмещённость и состоятельность — разные свойства. Оценка дисперсии ошибок
$\widehat{\sigma}^2 = \ESS/n$ смещена ($\E\widehat{\sigma}^2 =
\frac{n-k}{n}\sigma^2$), но состоятельна: смещение $-\frac{k}{n}\sigma^2$
исчезает при $n \to \infty$. Несмещённая оценка — $s^2 = \ESS/(n-k)$.
Наоборот, оценка $\widehat{\beta} = Y_1/X_1$ по одному наблюдению в модели
$Y_t = \beta X_t + \varepsilon_t$ несмещённая, но не состоятельна: она вообще
не использует остальные наблюдения, и её дисперсия не убывает с ростом $n$.
\end{remark}
\end{solution}

% -------------------------------------------------------------------
\problem{2}

Пусть есть модель $Y_t = \alpha + \beta_1X_{1t} + \beta_2X_{2t} +
\beta_3X_{3t} + \varepsilon_t$. Исследователь пытается проверить
мультиколлинеарность и оценивает вспомогательные регрессии: $X_1$ на $X_2,
X_3$ (коэффициент детерминации $R_1^2$), $X_2$ на $X_1, X_3$ ($R_2^2$), $X_3$ на
$X_1, X_2$ ($R_3^2$). Напишите, что такое мультиколлинеарность, и сделайте
вывод о её наличии или отсутствии в каждом варианте:
\begin{center}
\begin{tabular}{ccccc}
\toprule
 & Вариант 1 & Вариант 2 & Вариант 3 & Вариант 4 \\
\midrule
$R_1^2$ & 0,80 & 0,10 & 0,73 & 0,90 \\
$R_2^2$ & 0,25 & 0,65 & 0,60 & 0,65 \\
$R_3^2$ & 0,35 & 0,85 & 0,20 & 0,15 \\
\bottomrule
\end{tabular}
\end{center}

\begin{solution}
\textbf{Что такое мультиколлинеарность.} Это тесная линейная связь между
регрессорами. \emph{Строгая} мультиколлинеарность — точная линейная
зависимость столбцов матрицы $\mathbf{X}$: тогда
$\mathbf{X}^{\top}\mathbf{X}$ вырождена и МНК-оценки не существует.
\emph{Нестрогая} — сильная, но не точная связь: оценки существуют и остаются
несмещёнными, но их дисперсии велики. Признаки: высокий $R^2$ и значимый
$F$-тест при незначимых отдельных коэффициентах, неустойчивость оценок,
«неправильные» знаки.

Вспомогательный $R_j^2$ показывает, какая доля разброса $X_j$ объясняется
остальными регрессорами. По \eqref{eq:vif} переводим его в VIF —
во сколько раз из-за этой связи выросла дисперсия $\widehat{\beta}_j$:
\begin{center}
\begin{tabular}{ccccc}
\toprule
 & Вариант 1 & Вариант 2 & Вариант 3 & Вариант 4 \\
\midrule
$\mathrm{VIF}_1 = 1/(1 - R_1^2)$ & 5,00 & 1,11 & 3,70 & \textbf{10,00} \\
$\mathrm{VIF}_2 = 1/(1 - R_2^2)$ & 1,33 & 2,86 & 2,50 & 2,86 \\
$\mathrm{VIF}_3 = 1/(1 - R_3^2)$ & 1,54 & \textbf{6,67} & 1,25 & 1,18 \\
\bottomrule
\end{tabular}
\end{center}
Например, в варианте 1: $\mathrm{VIF}_1 = 1/(1 - 0{,}8) = 1/0{,}2 = 5$.

\textbf{Выводы} (порог $\mathrm{VIF} > 5$):
\begin{itemize}[leftmargin=2.2em, itemsep=2pt]
  \item \emph{Вариант 1.} $\mathrm{VIF}_1 = 5$ — ровно на пороге: $X_1$
        на 80\,\% объясняется остальными регрессорами. Умеренная
        мультиколлинеарность, связанная с $X_1$; по строгому порогу 10 её нет.
  \item \emph{Вариант 2.} $\mathrm{VIF}_3 = 6{,}67 > 5$ — мультиколлинеарность
        есть: $X_3$ тесно связан с $X_1$ и $X_2$. До серьёзной ($> 10$) не
        дотягивает.
  \item \emph{Вариант 3.} Все VIF меньше 5 — мультиколлинеарности нет.
  \item \emph{Вариант 4.} $\mathrm{VIF}_1 = 10$ — серьёзная
        мультиколлинеарность: дисперсия $\widehat{\beta}_1$ в 10 раз больше,
        чем была бы при регрессоре, не связанном с остальными, а стандартная
        ошибка — в $\sqrt{10} \approx 3{,}2$ раза.
\end{itemize}
\end{solution}

% -------------------------------------------------------------------
\problem{3}

\begin{enumerate}[label=\asbuk*), leftmargin=2.2em, itemsep=3pt]
  \item Какие есть способы решить проблему мультиколлинеарности?
  \item Напишите формулу скорректированного $R^2$. Объясните смысловое
        отличие от нескорректированного $R^2$.
  \item Сформулируйте байесовский информационный критерий (BIC): напишите
        формулу и объясните, зачем он используется.
  \item Сформулируйте информационный критерий Акаике (AIC): напишите формулу и
        объясните, зачем он используется.
\end{enumerate}

\begin{solution}
\textbf{а) Способы решения проблемы мультиколлинеарности.}
\begin{enumerate}[label=\arabic*), leftmargin=2.2em, itemsep=2pt]
  \item \emph{Исключить переменную (переменные) из модели.} Исключается одна
        из коррелирующих объясняющих переменных — та, у которой хуже
        показатели в однофакторном анализе (меньше $R^2$ парной регрессии $Y$
        на неё). Цена — смещение оставшихся коэффициентов, если исключённая
        переменная на самом деле влияет на $Y$.
  \item \emph{Получить дополнительные данные или новую выборку.} Например,
        заменить годовые данные квартальными: при большем числе наблюдений
        дисперсии оценок меньше, а связь регрессоров может оказаться слабее.
  \item \emph{Использовать предварительную информацию о параметрах.} Например,
        известное из теории ограничение на коэффициенты ($\beta_1 + \beta_2 =
        1$) или объединение коррелирующих объясняющих переменных в одну.
  \item \emph{Преобразовать данные.} Например, разделить на одну из
        коррелирующих объясняющих переменных (перейти к относительным
        величинам) или перейти к приростам.
\end{enumerate}

\medskip
\textbf{б) Скорректированный $R^2$.} По \eqref{eq:r2}
\[
  R^2 = 1 - \frac{\sum(y_t - \widehat{y}_t)^2}{\sum(y_t - \overline{y})^2},
  \qquad
  R^2_{\text{adj}} = 1 - \frac{\sum(y_t - \widehat{y}_t)^2/(n-k)}{\sum(y_t - \overline{y})^2/(n-1)}
                   = 1 - \frac{n-1}{n-k}\bigl(1 - R^2\bigr).
\]
Обычный $R^2$ показывает долю разброса $y$, объяснённую моделью, но не
убывает при добавлении \emph{любого} регрессора, даже бессмысленного:
остаточная сумма от этого может только уменьшиться. Поэтому сравнивать по
$R^2$ модели с разным числом регрессоров нельзя. В $R^2_{\text{adj}}$ каждая
сумма квадратов делится на своё число степеней свободы, то есть сравниваются
\emph{несмещённые оценки} дисперсии ошибок и дисперсии $y$. Добавление
регрессора увеличивает $R^2_{\text{adj}}$, только если выигрыш в остаточной
сумме перевешивает потерю степени свободы. Всегда $R^2_{\text{adj}} \le R^2$,
и $R^2_{\text{adj}}$ может быть отрицательным.

\medskip
\textbf{в) Критерий BIC (Шварца).} По \eqref{eq:ic}
\[
  \mathrm{BIC} = \frac{k\ln n}{n} + \ln\frac{\sum(y_t - \widehat{y}_t)^2}{n}
               = \frac{k\ln n}{n} + \ln\frac{\ESS}{n}.
\]

\textbf{г) Критерий AIC (Акаике).}
\[
  \mathrm{AIC} = \frac{2k}{n} + \ln\frac{\sum(y_t - \widehat{y}_t)^2}{n}
               = \frac{2k}{n} + \ln\frac{\ESS}{n}.
\]
В обоих случаях $n$ — число наблюдений, $k$ — число параметров (регрессоров,
включая константу).

\emph{Зачем.} Оба критерия используют для выбора между несколькими моделями,
оценёнными по одним и тем же данным, в том числе когда модели не вложены друг
в друга и $F$-тест на ограничения неприменим. Второе слагаемое — качество
подгонки (меньше остатки — меньше критерий), первое — штраф за число
параметров. Лучше модель с \emph{меньшим} значением критерия. Штраф BIC
$k\ln n/n$ больше штрафа AIC $2k/n$ при $\ln n > 2$, то есть уже при
$n \ge 8$, поэтому BIC строже наказывает за лишние параметры и выбирает более
экономные модели.
\end{solution}

% -------------------------------------------------------------------
\ifsolutions\clearpage\fi
\problem{4 (Магнус, упр.~2.22)}

В таблице приведены ежегодные значения денежной массы и национального дохода
некоторой гипотетической страны (в миллиардах денежных единиц).

\begin{center}
\begin{tabular}{ccc @{\qquad} ccc}
\toprule
Год & Денежная масса & Нац. доход & Год & Денежная масса & Нац. доход \\
\midrule
1981 & 2,0 & 5,0 & 1986 & 4,0 & 7,7 \\
1982 & 2,5 & 5,5 & 1987 & 4,2 & 8,4 \\
1983 & 3,2 & 6,0 & 1988 & 4,6 & 9,0 \\
1984 & 3,6 & 7,0 & 1989 & 4,8 & 9,7 \\
1985 & 3,3 & 7,2 & 1990 & 5,0 & 10,0 \\
\bottomrule
\end{tabular}
\end{center}

\begin{enumerate}[label=\asbuk*), leftmargin=2.2em, itemsep=3pt]
  \item Проведите регрессию национального дохода ($Y$) на денежную массу ($X$)
        и константу. Обоснуйте вывод основных формул.
  \item Рассчитайте $R^2$ модели.
  \item Постройте 95\,\%-ные доверительные интервалы для оцениваемых
        параметров. Можете ли вы отвергнуть гипотезу $\beta = 0$? $\beta = 1$?
\end{enumerate}

\begin{excel}
Введите данные двумя столбцами $X$ и $Y$. Суммы для ручного расчёта:
\xl{=СУММ(\dots)}, \xl{=СУММПРОИЗВ(X; X)}, \xl{=СУММПРОИЗВ(X; Y)}; суммы
квадратов отклонений — \xl{=КВАДРОТКЛ(\dots)}. Проверка: \ui{Данные} $\to$
\ui{Анализ данных} $\to$ \ui{Регрессия}, входной интервал $Y$ — национальный
доход, $X$ — денежная масса; галочка \ui{Уровень надёжности 95\,\%} даёт
доверительные интервалы.
\end{excel}

\begin{solution}
\textbf{а) Вывод формул МНК.} МНК выбирает $\alpha$ и $\beta$, при которых
сумма квадратов отклонений наблюдений от прямой минимальна:
\[
  S(\alpha, \beta) = \sum_{t=1}^{n}\bigl(Y_t - \alpha - \beta X_t\bigr)^2 \to \min_{\alpha,\,\beta}.
\]
Необходимое условие минимума — равенство нулю частных производных:
\[
  \frac{\partial S}{\partial\alpha} = -2\sum\bigl(Y_t - \alpha - \beta X_t\bigr) = 0,
  \qquad
  \frac{\partial S}{\partial\beta} = -2\sum X_t\bigl(Y_t - \alpha - \beta X_t\bigr) = 0 .
\]
Отсюда система нормальных уравнений:
\[
  \begin{cases}
    n\alpha + \beta\sum X_t = \sum Y_t, \\[2pt]
    \alpha\sum X_t + \beta\sum X_t^2 = \sum X_tY_t .
  \end{cases}
\]
Из первого уравнения (после деления на $n$) $\alpha = \overline{Y} -
\beta\overline{X}$: прямая проходит через точку средних. Подставляем во второе:
\[
  \bigl(\overline{Y} - \beta\overline{X}\bigr)\sum X_t + \beta\sum X_t^2 = \sum X_tY_t
  \quad\Longrightarrow
\]
\[
  \widehat{\beta} = \frac{\sum X_tY_t - n\overline{X}\,\overline{Y}}{\sum X_t^2 - n\overline{X}^2}
                  = \frac{\sum x_ty_t}{\sum x_t^2},
  \qquad
  \widehat{\alpha} = \overline{Y} - \widehat{\beta}\,\overline{X}.
\]
Это действительно минимум: $S$ — выпуклая квадратичная функция,
$\partial^2S/\partial\alpha^2 = 2n > 0$, а определитель матрицы вторых
производных $4\bigl(n\sum X_t^2 - (\sum X_t)^2\bigr) = 4n\sum x_t^2 > 0$, если
не все $X_t$ равны.

\medskip
\textbf{Расчёт.} $n = 10$. Суммы по таблице:
\[
  \sum X_t = 37{,}2,\quad \sum Y_t = 75{,}5,\quad \sum X_t^2 = 147{,}18,
\]
\[
  \sum Y_t^2 = 597{,}03,\quad \sum X_tY_t = 295{,}95 .
\]
Средние $\overline{X} = 3{,}72$, $\overline{Y} = 7{,}55$. Суммы в отклонениях:
\[
\begin{aligned}
  \sum x_t^2  &= \sum X_t^2 - \frac{(\sum X_t)^2}{n} = 147{,}18 - \frac{37{,}2^2}{10} = 8{,}796, \\
  \sum x_ty_t &= \sum X_tY_t - \frac{\sum X_t\sum Y_t}{n} = 295{,}95 - \frac{37{,}2\cdot 75{,}5}{10} = 15{,}09, \\
  \sum y_t^2  &= \sum Y_t^2 - \frac{(\sum Y_t)^2}{n} = 597{,}03 - \frac{75{,}5^2}{10} = 27{,}005 .
\end{aligned}
\]
По \eqref{eq:ols}
\[
  \widehat{\beta} = \frac{15{,}09}{8{,}796} = 1{,}716,
  \qquad
  \widehat{\alpha} = 7{,}55 - 1{,}716\cdot 3{,}72 = 1{,}168,
\]
\[
  \widehat{Y} = 1{,}168 + 1{,}716\,X .
\]
Рост денежной массы на 1 млрд сопровождается в среднем ростом национального
дохода на 1,72 млрд.

\medskip
\textbf{б) Коэффициент детерминации.} По \eqref{eq:ols}, \eqref{eq:r2}
\[
  \RSS = \widehat{\beta}^2\sum x_t^2 = 1{,}716^2\cdot 8{,}796 = 25{,}888,
  \qquad
  R^2 = \frac{\RSS}{\TSS} = \frac{25{,}888}{27{,}005} = 0{,}959 .
\]
Денежная масса объясняет 95,9\,\% разброса национального дохода.

\medskip
\textbf{в) Доверительные интервалы и гипотезы.} Остаточная сумма и оценка
дисперсии ошибок по \eqref{eq:s2}:
\[
  \ESS = \TSS - \RSS = 27{,}005 - 25{,}888 = 1{,}117,
  \qquad
  s^2 = \frac{\ESS}{n-2} = \frac{1{,}117}{8} = 0{,}1397 .
\]
Стандартные ошибки:
\[
  s_{\widehat{\beta}} = \sqrt{\frac{s^2}{\sum x_t^2}} = \sqrt{\frac{0{,}1397}{8{,}796}} = 0{,}126,
\]
\[
  s_{\widehat{\alpha}} = \sqrt{s^2\Bigl(\frac{1}{n} + \frac{\overline{X}^2}{\sum x_t^2}\Bigr)}
                       = \sqrt{0{,}1397\Bigl(\frac{1}{10} + \frac{3{,}72^2}{8{,}796}\Bigr)} = 0{,}483 .
\]
С $t_{0{,}025}(8) = 2{,}306$ по \eqref{eq:t}:
\[
  \beta\colon\ 1{,}716 \pm 2{,}306\cdot 0{,}126 = (1{,}425;\ 2{,}006),
  \qquad
  \alpha\colon\ 1{,}168 \pm 2{,}306\cdot 0{,}483 = (0{,}053;\ 2{,}283).
\]
Гипотезы о наклоне:
\[
\begin{aligned}
  H_0\colon \beta = 0\colon&\quad t = \frac{1{,}716 - 0}{0{,}126} = 13{,}6 > 2{,}306, \\[2pt]
  H_0\colon \beta = 1\colon&\quad t = \frac{1{,}716 - 1}{0{,}126} = 5{,}68 > 2{,}306 .
\end{aligned}
\]
Обе гипотезы отвергаются на 5\,\%-ном уровне. То же видно по
доверительному интервалу: ни 0, ни 1 в интервал $(1{,}425;\ 2{,}006)$ не
попадают.

\begin{remark}
Стандартная ошибка наклона — частный случай формулы \eqref{eq:se}: в парной
регрессии вспомогательная регрессия $X$ на константу имеет остаточную сумму
$\ESS_X = \sum x_t^2 = 8{,}796$, и
$s_{\widehat{\beta}} = s/\sqrt{\ESS_X} = \sqrt{0{,}1397}/\sqrt{8{,}796} = 0{,}126$.
\end{remark}
\end{solution}
